\documentclass{article}
\usepackage[utf8]{inputenc}
\usepackage{hyperref}
\usepackage{xcolor}
\usepackage{listings}
\usepackage{geometry}
\geometry{a4paper, margin=1in}

\title{\textbf{EMET - Quick Start Guide}}
\author{EMET Development Team}
\date{\today}

\begin{document}

\maketitle

\section*{Introduction}
EMET is a defensive-first vulnerability orchestration platform focused on evidence, provenance, and analyst workflows. This guide covers how to set up and run the EMET platform, either via Docker or manually.

\section{Docker Quick Start (Recommended)}

Using Docker Compose is the easiest way to start the EMET stack, which includes the frontend, backend, worker, and necessary databases.

\begin{enumerate}
    \item \textbf{Setup the Environment:}\\
    Copy the example environment variables file to create your own configuration.
    \begin{lstlisting}[language=bash, frame=single, backgroundcolor=\color{lightgray!20}]
cp .env.example .env
    \end{lstlisting}

    \item \textbf{Start the Stack:}\\
    Build and start all services in detached mode.
    \begin{lstlisting}[language=bash, frame=single, backgroundcolor=\color{lightgray!20}]
docker compose up --build -d
    \end{lstlisting}

    \item \textbf{Access the Application:}\\
    Once the containers are running, you can access the following services:
    \begin{itemize}
        \item \textbf{Frontend GUI:} \url{http://localhost:3000}
        \item \textbf{Backend Health Check:} \url{http://localhost:8000/api/health}
        \item \textbf{Backend Readiness:} \url{http://localhost:8000/api/readiness}
    \end{itemize}

    \item \textbf{Demo Credentials:}\\
    Use the following accounts to log into the frontend dashboard:
    \begin{itemize}
        \item Analyst: \texttt{analyst@emet.local} / Password: \texttt{emet}
        \item Guest (Read-Only): \texttt{guest} / Password: \texttt{emet}
    \end{itemize}
\end{enumerate}

\section{Manual Startup (Without Docker)}

If you prefer to run the components directly on your host machine, you will need to start the Frontend, Backend API, and Backend Worker separately.

\subsection{Frontend}
\begin{lstlisting}[language=bash, frame=single, backgroundcolor=\color{lightgray!20}]
cd frontend
npm install
npm run dev
\end{lstlisting}
The UI will be available at \url{http://localhost:3000}.

\subsection{Backend API}
\begin{lstlisting}[language=bash, frame=single, backgroundcolor=\color{lightgray!20}]
cd backend
python3 -m venv .venv
source .venv/bin/activate
pip install -r requirements.txt
uvicorn main:app --reload --host 0.0.0.0 --port 8000
\end{lstlisting}

\subsection{Backend Worker}
In a separate terminal window, start the task worker to process scan jobs.
\begin{lstlisting}[language=bash, frame=single, backgroundcolor=\color{lightgray!20}]
cd backend
source .venv/bin/activate
python worker.py
\end{lstlisting}

\section{Optional Services}

\textbf{Graph Mode (Neo4j)}\\
If you are testing Attack Paths and require Neo4j to be running:
\begin{lstlisting}[language=bash, frame=single, backgroundcolor=\color{lightgray!20}]
docker compose --profile graph up -d neo4j
\end{lstlisting}

\section*{Important Notes}
\begin{itemize}
    \item EMET is defensive-only and does not generate exploit code.
    \item Private/internal targets are blocked by default. To scan internal networks, you must set \texttt{ALLOW\_INTERNAL\_SCANNING=true} in your \texttt{.env} file.
    \item If AI model keys are missing, RAG capabilities and evaluation will continue to run in an explicit fallback mode.
\end{itemize}

\end{document}
